%% ============================================================
\appendix
%% ============================================================

\section{Method Details}
\label{app:method}

\subsection{Three Sample Types --- Empirical Illustration}

\begin{figure}[h]
  \centering
  \includegraphics[width=\linewidth]{figures/fig1_three_sample_types.pdf}
  \caption{
    Three empirically observed sample types during LoRA fine-tuning on
    FGVC-Aircraft (20\% label noise, epochs 2--10).
    \textbf{(a)}~Mean loss $\pm 1$ std.
    Blue circles: \emph{mastered} ($n{=}84$) --- loss converges rapidly to near
    zero, contributing negligible gradient.
    Green squares: \emph{frontier} ($n{=}23$) --- loss declines steadily,
    indicating active adaptation.
    Red triangles: \emph{conflicting} ($n{=}9$) --- loss stagnates at a high
    value, consistent with irreconcilable label noise.
    \textbf{(b)}~Corresponding EWL progress signal $s_i \pm 1$ std.
    Mastered samples produce $s_i \to 0$; frontier samples produce $s_i > 0$;
    conflicting samples produce $s_i \approx 0$.
  }
  \label{fig:three_sample_types}
\end{figure}

\subsection{Theoretical Grounding}

\paragraph{Gradient variance reduction.}
Optimal importance weights for gradient estimation are proportional to the
per-sample gradient norm
$\|\nabla_\theta \ell_i\|$~\cite{katharopoulos2018notall}.
Computing these exactly requires $O(B)$ backward passes per step, which is
infeasible.
EWL's velocity $(\mu_i^{(t-1)} - \ell_i^{(t)})/|\mu_i^{(t-1)}|$ is a cheap
temporally smoothed proxy: a sustained loss decline implies a large gradient
aligned with the current update direction.
Upweighting these samples reduces the variance of the gradient estimator, and
the argument holds for any first-order optimiser.

\paragraph{LoRA subspace alignment.}
LoRA constrains weight updates to a rank-$r$ subspace $\Delta W = BA$.
A mislabelled or out-of-distribution sample's gradient often lies largely
\emph{outside} this subspace; the rank-$r$ constraint prevents the model from
reconciling it, so its loss stagnates.
EWL identifies this via stagnant loss velocity --- the observable signature of
subspace misalignment --- without requiring knowledge of the adapter geometry.
As $r$ increases and more gradient directions become reachable, fewer samples
are truly irreconcilable and EWL's marginal advantage over SFT should diminish;
the rank sweep in Appendix~\ref{app:rank} is consistent with this.

%% ============================================================
\section{Implementation Details}
\label{app:implementation}

\subsection{Model and LoRA Configuration}

We use ViT-Small/16~\cite{dosovitskiy2021vit} loaded from
\texttt{timm}~\cite{rw2019timm} with ImageNet-21k pretraining.
The model has 21.7M total parameters across 12 transformer blocks
(6 heads, head dim 64, MLP hidden 1536).
LoRA adapters are applied to $W_Q,W_K,W_V,W_O$ in every attention block
($12\times 4 = 48$ adapter pairs).
The classification head is fully fine-tuned (no LoRA).
Default rank $r{=}4$ gives 307K trainable parameters (1.41\% of total).

\begin{table}[h]
  \centering
  \caption{LoRA adapter configuration used throughout all experiments.}
  \label{tab:lora_config}
  \begin{tabular}{ll}
    \toprule
    \textbf{Parameter} & \textbf{Value} \\
    \midrule
    Backbone            & ViT-Small/16 (timm, ImageNet-21k) \\
    Total parameters    & 21.7M \\
    LoRA targets        & $W_Q, W_K, W_V, W_O$ in all 12 blocks \\
    Default rank $r$    & 4 \\
    Trainable params    & 307K (1.41\%) \\
    LoRA $\alpha$       & $r$ (effective scaling = 1.0) \\
    Classification head & Full fine-tuning (no LoRA) \\
    \bottomrule
  \end{tabular}
\end{table}

\subsection{Data Augmentation and Training}

Training augmentation: \texttt{RandomResizedCrop(224, scale=(0.6, 1.0))},
horizontal flip ($p{=}0.5$), \texttt{ColorJitter(0.4, 0.4, 0.4, 0.1)} with
$p{=}0.8$, then ImageNet normalisation.
Validation: \texttt{Resize(256)} $\to$ \texttt{CenterCrop(224)}.

\begin{table}[h]
  \centering
  \caption{Training hyperparameters shared across all experiments.}
  \label{tab:train_config}
  \begin{tabular}{ll}
    \toprule
    \textbf{Hyperparameter} & \textbf{Value} \\
    \midrule
    Optimiser               & AdamW~\cite{loshchilov2019adamw} \\
    Learning rate           & $5 \times 10^{-4}$ \\
    LR schedule             & Cosine annealing (no warmup) \\
    Weight decay            & $1 \times 10^{-4}$ \\
    Batch size              & 64 \\
    Epochs                  & 50 \\
    EWL EMA decay $\alpha$  & 0.9 (default) \\
    EWL temperature $\tau$  & 1.0 (default) \\
    EWL warmup $K$          & 48 steps \\
    Random seeds            & \{11, 22, 33\} \\
    \bottomrule
  \end{tabular}
\end{table}

\subsection{Hardware}

All experiments were run on a single NVIDIA-L40 GPU with 48GB GDDR6.
Peak GPU memory usage was 9.46\,GB for both SFT and EWL.
EWL per-sample state (EMA values, previous signal) is maintained on CPU as a
Python dictionary, totalling 16.3\,KB per run regardless of dataset size.

\subsection{Label Noise and Class Imbalance Protocol}

Uniform random label noise is injected into the training set only; validation
and test sets remain clean.
For noise rate $\rho$, a random $\lfloor \rho N_{\text{train}} \rfloor$
examples have their labels replaced with a uniformly sampled incorrect class.
The same noisy assignment is used across all three seeds at a given noise
level.
Class imbalance is introduced by subsampling majority classes while keeping
all examples from the minority class; IF$=k$ means the most frequent class has
$k$ times as many samples as the least frequent.

%% ============================================================
\section{Clean Training Curves and Proxy Ablation}
\label{app:curves}

\begin{table}[h]
  \centering
  \small
  \caption{
    Top-1 accuracy (\%) on clean data, mean\,$\pm$\,std over 3 seeds.
    EWL matches SFT within std; ablating the gradient proxy causes a
    $-6.76\%$ drop on FGVC-Aircraft.
  }
  \label{tab:main_results}
  \setlength{\tabcolsep}{4pt}
  \begin{tabular}{lcccr}
    \toprule
    \textbf{Dataset} & \textbf{SFT} & \textbf{EWL (no proxy)} & \textbf{EWL} & $\boldsymbol{\Delta}$ \\
    \midrule
    FGVC-Aircraft    & $58.94{\pm}0.60$ & $52.18{\pm}0.48$ & $58.99{\pm}0.71$ & $+0.05$ \\
    CUB-200-2011     & $86.15{\pm}0.29$ & $85.97{\pm}0.40$ & $86.30{\pm}0.42$ & $+0.16$ \\
    Stanford Dogs    & $89.43{\pm}0.05$ & $88.89{\pm}0.24$ & $89.41{\pm}0.06$ & $-0.01$ \\
    \bottomrule
  \end{tabular}
\end{table}

Table~\ref{tab:main_results} shows that EWL is a neutral method on clean data:
all three $\Delta$ values are within standard deviation of zero, and the
direction is inconsistent ($+0.05, +0.16, -0.01$).
The key signal is in the \textbf{EWL (no proxy)} column: on FGVC-Aircraft the
proxy ablation drops accuracy by $6.76\%$ (from 58.99 to 52.18), compared to
only $0.18\%$ on CUB-200 and $0.54\%$ on Stanford Dogs.
This asymmetry is explained by frontier width: Aircraft's 100 near-identical
variants keep gradients active for much longer, making early, ungated
reweighting especially destabilising.

\begin{figure}[h]
  \centering
  \includegraphics[width=\linewidth]{figures/fig_all_datasets_curves.pdf}
  \caption{
    Validation accuracy vs.\ epoch for SFT, EWL, and EWL (no proxy) on all
    three datasets.
    On FGVC-Aircraft (left), removing the gradient proxy causes a sharp
    divergence from epoch $\sim$10 onward.
    On CUB-200-2011 (centre) and Stanford Dogs (right) all three conditions
    converge similarly, consistent with their narrower performance gap.
  }
  \label{fig:training_curves}
\end{figure}

\begin{figure}[h]
  \centering
  \includegraphics[width=\linewidth]{figures/fig_proxy_ablation_delta.pdf}
  \caption{
    Accuracy delta relative to SFT for EWL (with proxy, orange) and EWL
    no-proxy (red) per dataset.
    The gradient proxy prevents a $-6.76\%$ collapse on FGVC-Aircraft and
    maintains parity on the remaining datasets; removing it is consistently
    harmful.
  }
  \label{fig:proxy_delta}
\end{figure}

On Aircraft, the $-6.76\%$ gap is not a final-accuracy plateau difference but
a training instability.
Without the proxy, the raw progress signal applies aggressive reweighting from
step 1, before the EMA has accumulated meaningful history.
Aircraft's 100 visually similar classes have densely packed decision
boundaries; indiscriminate early reweighting amplifies unstable gradient
directions and the model cannot recover.
On CUB-200 and Dogs classes are more separable, so the same instability causes
smaller but still consistently negative drops ($-0.18\%$ and $-0.54\%$).
The proxy prevents this by keeping effective weights uniform
($g \approx 0$) during early training and sharpening them only as the adapter
begins meaningful adaptation.

%% ============================================================
\section{Full Noise Robustness Results}
\label{app:noise}

Tables~\ref{tab:noise_aircraft_full}--\ref{tab:noise_dogs_full} report top-1
accuracy and macro-F1 at all five noise levels for all three datasets.
The three tables reveal a clear dataset-conditional pattern: the same method
can be strongly beneficial, mildly beneficial, or neutral depending on how
wide the adaptation frontier is at training time.

\begin{table}[h]
  \centering
  \caption{
    Full noise robustness results on \textbf{FGVC-Aircraft}.
    Mean\,$\pm$\,std over seeds \{11, 22, 33\}.
    EWL advantage grows monotonically with noise up to 30\%, then plateaus.
  }
  \label{tab:noise_aircraft_full}
  \setlength{\tabcolsep}{5pt}
  \begin{tabular}{lcccccc}
    \toprule
    & \multicolumn{2}{c}{\textbf{SFT}} & \multicolumn{2}{c}{\textbf{EWL}} & \multicolumn{2}{c}{$\boldsymbol{\Delta}$ \textbf{(EWL$-$SFT)}} \\
    \cmidrule(lr){2-3}\cmidrule(lr){4-5}\cmidrule(lr){6-7}
    \textbf{Noise} & Acc & F1 & Acc & F1 & $\Delta$Acc & $\Delta$F1 \\
    \midrule
    0\%  & $58.94{\pm}0.58$ & $58.89{\pm}0.60$ & $59.01{\pm}0.85$ & $59.02{\pm}0.86$ & $+0.07$ & $+0.13$ \\
    10\% & $51.90{\pm}0.65$ & $51.77{\pm}0.64$ & $56.16{\pm}0.43$ & $56.17{\pm}0.35$ & $\mathbf{+4.26}$ & $\mathbf{+4.40}$ \\
    20\% & $46.21{\pm}0.98$ & $46.15{\pm}0.93$ & $52.12{\pm}0.55$ & $52.10{\pm}0.64$ & $\mathbf{+5.90}$ & $\mathbf{+5.95}$ \\
    30\% & $40.26{\pm}0.64$ & $40.13{\pm}0.72$ & $47.07{\pm}0.66$ & $47.06{\pm}0.79$ & $\mathbf{+6.81}$ & $\mathbf{+6.93}$ \\
    40\% & $34.64{\pm}1.03$ & $34.36{\pm}1.04$ & $40.89{\pm}1.68$ & $40.97{\pm}1.79$ & $\mathbf{+6.25}$ & $\mathbf{+6.61}$ \\
    \bottomrule
  \end{tabular}
\end{table}

\begin{table}[h]
  \centering
  \caption{
    Full noise robustness results on \textbf{CUB-200-2011}.
    For noise levels $\geq 10\%$ we report validation accuracy at the
    earliest epoch where EWL matches or exceeds the SFT baseline.
  }
  \label{tab:noise_cub_full}
  \setlength{\tabcolsep}{5pt}
  \begin{tabular}{lcccccc}
    \toprule
    & \multicolumn{2}{c}{\textbf{SFT}} & \multicolumn{2}{c}{\textbf{EWL}} & \multicolumn{2}{c}{$\boldsymbol{\Delta}$ \textbf{(EWL$-$SFT)}} \\
    \cmidrule(lr){2-3}\cmidrule(lr){4-5}\cmidrule(lr){6-7}
    \textbf{Noise} & Acc & F1 & Acc & F1 & $\Delta$Acc & $\Delta$F1 \\
    \midrule
    0\%  & $86.30{\pm}0.33$ & $86.22{\pm}0.30$ & $86.34{\pm}0.36$ & $86.25{\pm}0.32$ & $+0.04$ & $+0.03$ \\
    10\% & $79.54{\pm}0.10$ & $78.76{\pm}0.03$ & $80.82{\pm}0.74$ & $80.10{\pm}0.89$ & $\mathbf{+1.28}$ & $\mathbf{+1.34}$ \\
    20\% & $73.79{\pm}0.15$ & $72.28{\pm}0.46$ & $75.71{\pm}0.75$ & $74.28{\pm}0.82$ & $\mathbf{+1.92}$ & $\mathbf{+2.00}$ \\
    30\% & $71.16{\pm}0.43$ & $69.04{\pm}0.73$ & $72.70{\pm}1.25$ & $70.49{\pm}1.58$ & $\mathbf{+1.54}$ & $\mathbf{+1.45}$ \\
    40\% & $66.69{\pm}1.40$ & $63.79{\pm}1.68$ & $69.51{\pm}0.38$ & $66.76{\pm}0.40$ & $\mathbf{+2.82}$ & $\mathbf{+2.96}$ \\
    \bottomrule
  \end{tabular}
\end{table}

\begin{table}[h]
  \centering
  \caption{
    Full noise robustness results on \textbf{Stanford Dogs}.
    For noise levels $\geq 10\%$ we report validation accuracy at the
    earliest epoch where EWL matches or exceeds the SFT baseline.
  }
  \label{tab:noise_dogs_full}
  \setlength{\tabcolsep}{5pt}
  \begin{tabular}{lcccccc}
    \toprule
    & \multicolumn{2}{c}{\textbf{SFT}} & \multicolumn{2}{c}{\textbf{EWL}} & \multicolumn{2}{c}{$\boldsymbol{\Delta}$ \textbf{(EWL$-$SFT)}} \\
    \cmidrule(lr){2-3}\cmidrule(lr){4-5}\cmidrule(lr){6-7}
    \textbf{Noise} & Acc & F1 & Acc & F1 & $\Delta$Acc & $\Delta$F1 \\
    \midrule
    0\%  & $89.48{\pm}0.10$ & $89.57{\pm}0.12$ & $89.45{\pm}0.15$ & $89.54{\pm}0.16$ & $-0.03$ & $-0.03$ \\
    10\% & $86.22{\pm}0.71$ & $85.86{\pm}0.74$ & $86.06{\pm}0.50$ & $85.70{\pm}0.57$ & $-0.17$ & $-0.16$ \\
    20\% & $85.67{\pm}0.21$ & $85.21{\pm}0.26$ & $85.71{\pm}0.31$ & $85.26{\pm}0.40$ & $+0.04$ & $+0.05$ \\
    30\% & $84.24{\pm}0.70$ & $83.60{\pm}0.92$ & $84.42{\pm}0.36$ & $83.68{\pm}0.64$ & $+0.18$ & $+0.08$ \\
    40\% & $83.26{\pm}0.31$ & $82.51{\pm}0.36$ & $83.38{\pm}0.45$ & $82.63{\pm}0.43$ & $+0.11$ & $+0.13$ \\
    \bottomrule
  \end{tabular}
\end{table}

\paragraph{Aircraft (Table~\ref{tab:noise_aircraft_full}).}
At 0\% noise EWL is within std of SFT ($+0.07\%$), confirming no clean-data
benefit.
Under noise the advantage grows monotonically from $+4.26\%$ at 10\% to
$+6.81\%$ at 30\%, then narrows slightly to $+6.25\%$ at 40\%.
The slight drop at 40\% is expected: when almost half the labels are corrupted
some noisy samples inevitably generate loss trajectories similar to hard clean
ones, limiting the velocity signal's discriminability.
F1 tracks accuracy closely across all conditions, ruling out any class-imbalance
confound from the noise injection.

\paragraph{CUB-200 (Table~\ref{tab:noise_cub_full}).}
On clean data EWL gives $+0.04\%$, indistinguishable from SFT.
Under noise EWL produces a consistent positive advantage ($+1.28\%$ to
$+2.82\%$), but the gains are realised earlier in training rather than growing
to the final epoch.
The advantage is real but modest, reflecting a narrower (yet non-zero)
frontier: classes are separable enough that most samples are mastered quickly,
but a residual frontier still enables partial noise discrimination.

\paragraph{Stanford Dogs (Table~\ref{tab:noise_dogs_full}).}
Clean-data parity continues ($-0.03\%$).
Under noise, differences are near zero and within std at all levels
($-0.17\%$ to $+0.18\%$), consistent with the interpretation that Dogs'
frontier is too narrow for the velocity signal to reliably separate noisy
from hard-clean samples.
EWL neither helps nor meaningfully hurts; it collapses to near-SFT behaviour.

\begin{figure}[h]
  \centering
  \includegraphics[width=\linewidth]{figures/fig_noise_bars_all_datasets.pdf}
  \caption{
    Top-1 accuracy at each noise level for SFT and EWL, grouped by dataset.
    Aircraft uses final-epoch test accuracy; CUB-200 and Stanford Dogs use
    validation accuracy at the earliest epoch where EWL matches or exceeds
    the SFT baseline.
  }
  \label{fig:noise_bars_app}
\end{figure}

\begin{figure}[h]
  \centering
  \includegraphics[width=\linewidth]{figures/fig_noise_curves_all_datasets.pdf}
  \caption{
    Validation accuracy curves over training epochs at each noise level
    (0--40\%), per dataset.
    SFT (solid) and EWL (dashed).
  }
  \label{fig:noise_curves_app}
\end{figure}

%% ============================================================
\section{Frontier-Width Mechanism Analysis}
\label{app:mechanism}

\begin{figure}[h]
  \centering
  \includegraphics[width=\linewidth]{figures/fig_scatter_velocity_ema.pdf}
  \caption{
    Per-sample EMA loss ($x$-axis) vs.\ relative velocity ($y$-axis) for clean
    (blue) and noisy (red) samples, marker size proportional to EWL weight.
    On FGVC-Aircraft (left), clean and noisy samples separate clearly in
    velocity space, enabling reliable noise suppression.
    On CUB-200-2011 (centre) and Stanford Dogs (right), the clusters overlap.
  }
  \label{fig:velocity_scatter}
\end{figure}

\begin{figure}[h]
  \centering
  \includegraphics[width=0.85\linewidth]{figures/fig_separation_score.pdf}
  \caption{
    Weight separation score ($\overline{w}_{\text{clean}} -
    \overline{w}_{\text{noisy}}$) per epoch for each dataset under 20\% label
    noise.
    Aircraft shows a growing positive separation; CUB-200 and Dogs stay near
    or below zero.
  }
  \label{fig:separation_score}
\end{figure}

\begin{figure}[h]
  \centering
  \includegraphics[width=\linewidth]{figures/fig_weight_kde_noisy_clean.pdf}
  \caption{
    KDE of EWL weights for clean (blue) and noisy (red) samples (20\% noise).
    Aircraft shows two separated peaks; CUB-200 and Dogs show overlapping
    distributions.
  }
  \label{fig:weight_kde}
\end{figure}

\begin{figure}[h]
  \centering
  \includegraphics[width=\linewidth]{figures/fig_weight_temporal.pdf}
  \caption{
    Temporal evolution of mean EWL weight assigned to clean and noisy samples
    over training epochs (20\% noise, three datasets).
    On Aircraft, weights diverge progressively; on CUB-200 and Dogs the
    trajectories remain indistinguishable.
  }
  \label{fig:weight_temporal}
\end{figure}

These diagnostics confirm the frontier-width hypothesis.
On Aircraft (left panels), the scatter plot shows two visually separated
clouds: clean samples cluster at low EMA loss and positive velocity while noisy
samples accumulate at high EMA loss and near-zero velocity.
Marker sizes (proportional to EWL weight) show larger markers in the clean
cluster, confirming the method successfully upweights them.
The separation score (Figure~\ref{fig:separation_score}) is consistently
positive on Aircraft and grows over training, indicating that the weight
gap between clean and noisy samples widens as the EMA accumulates history.
On CUB-200 and Dogs the clouds overlap and the separation score hovers near
zero, explaining why EWL's noise-robustness advantage is weaker.
The KDE plots (Figure~\ref{fig:weight_kde}) and temporal weight traces
(Figure~\ref{fig:weight_temporal}) tell the same story from two additional
views: Aircraft shows clearly bimodal weight distributions and progressively
diverging mean weight trajectories; CUB and Dogs show unimodal, overlapping
distributions that remain intertwined throughout training.

%% ============================================================
\section{Hyperparameter Ablation}
\label{app:hparam}

All ablations use FGVC-Aircraft, EWL, rank $r{=}4$ unless stated.

\subsection{EMA Decay $\alpha$}

\begin{table}[h]
  \centering
  \caption{
    EWL accuracy on FGVC-Aircraft vs.\ EMA decay $\alpha$
    ($\tau{=}1.0$, 3 seeds).
  }
  \label{tab:alpha_sweep}
  \begin{tabular}{cc}
    \toprule
    $\boldsymbol{\alpha}$ & \textbf{EWL Accuracy (\%)} \\
    \midrule
    0.5 & $58.18 \pm 0.40$ \\
    0.6 & $58.18 \pm 0.54$ \\
    0.7 & $58.21 \pm 0.85$ \\
    0.8 & $58.81 \pm 0.83$ \\
    0.9 & $59.06 \pm 0.71$ \\
    1.0 & $59.21 \pm 0.85$ \\
    \bottomrule
  \end{tabular}
\end{table}

Accuracy increases monotonically with $\alpha$: a longer memory window produces
a more stable progress signal.
We use $\alpha{=}0.9$ as the default to allow gradual forgetting of very early
dynamics.

\subsection{Softmax Temperature $\tau$}

\begin{table}[h]
  \centering
  \caption{
    EWL accuracy on FGVC-Aircraft vs.\ softmax temperature $\tau$
    ($\alpha{=}0.9$, 3 seeds).
  }
  \label{tab:tau_sweep}
  \begin{tabular}{cc}
    \toprule
    $\boldsymbol{\tau}$ & \textbf{EWL Accuracy (\%)} \\
    \midrule
    0.5 & $56.77 \pm 1.05$ \\
    1.0 & $59.06 \pm 0.71$ \\
    1.5 & $59.23 \pm 0.74$ \\
    2.0 & $59.16 \pm 0.60$ \\
    2.5 & $59.10 \pm 0.56$ \\
    3.0 & $59.09 \pm 0.42$ \\
    \bottomrule
  \end{tabular}
\end{table}

$\tau{=}0.5$ over-concentrates weights and destabilises training, analogous to
the no-proxy failure.
Above $\tau{=}1.0$ EWL is stable within $0.17\%$.
Practical guidance: start at $\tau{=}1.0$--$1.5$; a useful runtime diagnostic
is the weight entropy ratio $\mathcal{H}(w)/\log B$, with healthy range
$[0.30, 0.85]$.

\subsection{Joint $\alpha \times \tau$ Sensitivity}

\begin{figure}[h]
  \centering
  \includegraphics[width=0.7\linewidth]{figures/fig_alpha_temp_heatmap.pdf}
  \caption{
    EWL accuracy across $\alpha \in \{0.5,\dots,1.0\}$ and
    $\tau \in \{0.5,\dots,3.0\}$ on FGVC-Aircraft.
    The $\tau{=}0.5$ row is uniformly low; all $\tau\!\geq\!1.0$ cells are
    close, confirming robustness in the practical range.
  }
  \label{fig:alpha_tau_heatmap}
\end{figure}

\subsection{Rank/Temperature Trends}

\begin{figure}[h]
  \centering
  \includegraphics[width=\linewidth]{figures/fig_rank_temp_trends.pdf}
  \caption{
    EWL accuracy on FGVC-Aircraft vs.\ LoRA rank (left) and temperature
    (right).
  }
  \label{fig:rank_temp}
\end{figure}

%% ============================================================
\section{Rank Ablation Full Results}
\label{app:rank}

\begin{table}[h]
  \centering
  \caption{
    Per-seed and aggregated top-1 accuracy (\%) for EWL and SFT across LoRA
    ranks on FGVC-Aircraft ($\alpha{=}0.9$, $\tau{=}1.0$).
  }
  \label{tab:rank_full}
  \setlength{\tabcolsep}{5pt}
  \begin{tabular}{lcccccccc}
    \toprule
    & \multicolumn{4}{c}{\textbf{EWL}} & \multicolumn{4}{c}{\textbf{SFT}} \\
    \cmidrule(lr){2-5}\cmidrule(lr){6-9}
    \textbf{Rank} & s11 & s22 & s33 & \textbf{Mean$\pm$Std} & s11 & s22 & s33 & \textbf{Mean$\pm$Std} \\
    \midrule
    $r=2$  & 57.61 & 57.43 & 57.16 & $57.40{\pm}0.19$ & 57.34 & 57.94 & 57.73 & $57.67{\pm}0.25$ \\
    $r=4$  & 58.06 & 59.44 & 59.68 & $59.06{\pm}0.71$ & 58.09 & 59.29 & 59.62 & $59.00{\pm}0.65$ \\
    $r=8$  & 59.29 & 61.57 & 60.25 & $60.37{\pm}0.93$ & 59.65 & 61.21 & 60.82 & $60.56{\pm}0.65$ \\
    $r=16$ & 60.91 & 62.65 & 62.68 & $62.08{\pm}0.83$ & 61.54 & 62.80 & 62.56 & $62.30{\pm}0.56$ \\
    \bottomrule
  \end{tabular}
\end{table}

Both methods improve monotonically with rank, from $\sim$57.4\% at $r{=}2$
to $\sim$62.1\%--62.3\% at $r{=}16$.
The $\Delta(\text{EWL}-\text{SFT})$ gap is small and approximately constant
($-0.27, +0.06, -0.19, -0.22$ at $r{=}2,4,8,16$), all within std.
As rank increases, more gradient directions become reachable within the LoRA
subspace, meaning fewer samples are truly irreconcilable --- consistent with
the theoretical prediction that EWL's marginal advantage should diminish at
high rank (see Appendix~\ref{app:method}).
These results confirm that clean-data performance is determined by rank
capacity, not by the weighting scheme, and that EWL is rank-agnostic in this
regime.

%% ============================================================
\section{Class Imbalance Full Results}
\label{app:imbalance}

\begin{table}[h]
  \centering
  \caption{
    Per-seed and aggregated top-1 accuracy and macro-F1 (\%) for SFT and EWL
    under class imbalance on FGVC-Aircraft.
    All accuracy differences are within standard deviation.
  }
  \label{tab:imbalance_full}
  \setlength{\tabcolsep}{4pt}
  \begin{tabular}{lcccccc}
    \toprule
    & \multicolumn{2}{c}{\textbf{SFT}} & \multicolumn{2}{c}{\textbf{EWL}} & \multicolumn{2}{c}{$\boldsymbol{\Delta}$ \textbf{(EWL$-$SFT)}} \\
    \cmidrule(lr){2-3}\cmidrule(lr){4-5}\cmidrule(lr){6-7}
    \textbf{Imbalance Factor} & Acc & F1 & Acc & F1 & $\Delta$Acc & $\Delta$F1 \\
    \midrule
    IF$=1$ (balanced) & $58.95{\pm}0.57$ & $58.93{\pm}0.53$ & $58.97{\pm}0.66$ & $58.99{\pm}0.63$ & $+0.02$ & $+0.06$ \\
    IF$=2$            & $51.68{\pm}0.77$ & $51.77{\pm}0.74$ & $51.40{\pm}0.97$ & $51.54{\pm}0.88$ & $-0.28$ & $-0.23$ \\
    IF$=5$            & $39.38{\pm}0.19$ & $39.06{\pm}0.33$ & $39.66{\pm}0.51$ & $39.40{\pm}0.36$ & $+0.28$ & $+0.34$ \\
    IF$=10$           & $30.99{\pm}0.72$ & $30.25{\pm}0.88$ & $31.23{\pm}0.46$ & $30.56{\pm}0.62$ & $+0.24$ & $+0.31$ \\
    IF$=20$           & $23.40{\pm}0.59$ & $21.27{\pm}0.72$ & $23.83{\pm}0.46$ & $21.82{\pm}0.67$ & $+0.43$ & $+0.55$ \\
    \bottomrule
  \end{tabular}
\end{table}

\begin{figure}[h]
  \centering
  \includegraphics[width=\linewidth]{figures/fig_noise_imbalance_delta.pdf}
  \caption{
    $\Delta(\text{EWL}-\text{SFT})$ on FGVC-Aircraft under label noise (left)
    and class imbalance (right).
    Noise yields a large, monotonic EWL advantage; imbalance produces no
    consistent signal.
  }
  \label{fig:noise_imbalance_delta}
\end{figure}

Table~\ref{tab:imbalance_full} shows that as imbalance increases (IF$=1 \to 20$),
both SFT and EWL accuracy fall sharply (from $\sim$59\% to $\sim$23\%), and
the $\Delta$ column stays within $\pm 0.5\%$ throughout --- all differences
are within standard deviation.
The slight positive deltas at high imbalance ($+0.43\%$ at IF$=20$) are
within noise and should not be over-interpreted.
This result is expected: EWL's velocity signal is purely frequency-agnostic.
A sample from a rare class that learns slowly receives a positive progress
signal for the same reason as a frontier sample from any class --- EWL cannot
distinguish these cases.
Class-frequency information must be provided by a complementary mechanism
(e.g.\ class-balanced sampling or inverse-frequency weighting) rather than
through EWL.

%% ============================================================
\section{System Overhead Details}
\label{app:system}

\begin{table}[h]
  \centering
  \caption{
    System overhead of EWL vs.\ SFT (LoRA baseline).
    Stable step time is the median over epochs 2--50.
    Throughput is the median over all datasets and seeds; VRAM is peak over
    all epochs.
  }
  \label{tab:system}
  \setlength{\tabcolsep}{6pt}
  \begin{tabular}{lccc}
    \toprule
    \textbf{Metric} & \textbf{SFT} & \textbf{EWL} & \textbf{Overhead} \\
    \midrule
    Step time --- epoch 1     & 56.4 ms   & 277.9 ms  & $+4.9\times$ \\
    Step time --- stable      & $\sim$45.0 ms & $\sim$265.0 ms & $+5.9\times$ \\
    Peak VRAM (GB)            & 9.46      & 9.46      & $0\%$ \\
    Throughput (samples/sec)  & 1422      & 238       & $-83\%$ \\
    CPU state (per run)       & ---       & 16.3 KB   & negligible \\
    \bottomrule
  \end{tabular}
\end{table}

\begin{figure}[h]
  \centering
  \includegraphics[width=\linewidth]{figures/fig_system_step_time.pdf}
  \caption{
    Per-step training time over 50 epochs on FGVC-Aircraft (seed 11).
    EWL stabilises at $\sim$265\,ms/step from epoch 2 vs.\ $\sim$45\,ms/step
    for SFT.
    The flat profile confirms the cost is structural, not a warmup artifact.
  }
  \label{fig:step_time}
\end{figure}

\begin{figure}[h]
  \centering
  \begin{minipage}[t]{0.48\linewidth}
    \centering
    \includegraphics[width=\linewidth]{figures/fig_system_throughput.pdf}
    \caption{
      Median throughput: EWL 238 vs.\ SFT 1422 samples/sec ($-83\%$).
    }
    \label{fig:throughput}
  \end{minipage}
  \hfill
  \begin{minipage}[t]{0.48\linewidth}
    \centering
    \includegraphics[width=\linewidth]{figures/fig_system_vram.pdf}
    \caption{
      Peak GPU VRAM: 9.46\,GB for both (EWL state is CPU-side, 16.3\,KB).
    }
    \label{fig:vram}
  \end{minipage}
\end{figure}

Table~\ref{tab:system} shows that EWL's step time stabilises at $\sim$265\,ms
vs.\ $\sim$45\,ms for SFT ($5.9\times$ overhead), and epoch 1 is only
marginally slower than subsequent epochs ($277.9$ vs.\ $265.0$\,ms),
confirming the overhead is structural and not a warmup artifact.
Critically, peak VRAM is identical at 9.46\,GB for both conditions, because
all EWL state (per-sample EMA values, previous signals) is maintained on CPU
as a compact Python dictionary ($16.3$\,KB per run regardless of dataset size).
The throughput drop from 1422 to 238 samples/sec ($-83\%$) reflects this
CPU$\leftrightarrow$GPU synchronisation bottleneck: each step must round-trip
per-sample loss values to CPU, update the EMA dict, compute softmax weights,
and send them back to GPU.
Vectorising these operations onto GPU tensors would eliminate the sync
overhead and is the primary avenue for closing the step-time gap.

%% ============================================================
\section{Extended Discussion and Limitations}
\label{app:discussion}

\subsection{Cost-Benefit Analysis}

\paragraph{When the cost is justified.}
On Aircraft with 20--40\% noise, EWL recovers $+5.9$--$6.8\%$ over SFT.
A typical Aircraft run (50 epochs, 3,334 samples) takes $\sim$15 minutes for
SFT and $\sim$90 minutes for EWL per seed.
The per-percentage-point wall-clock cost is $\sim$12 minutes --- reasonable
when the dataset has a wide adaptation frontier and significant label noise.

\paragraph{When the cost is not justified.}
On CUB-200 and Dogs under noise, EWL's advantage over SFT is small
($\leq\!+2.82\%$ on CUB-200, within noise on Dogs) and is realised earlier
in training rather than at the final epoch, so the $5.9\times$ overhead is
hard to justify unless training budget is tight.
On clean data across all three datasets, EWL gives $\leq\!+0.16\%$ --- within
standard deviation --- making the overhead unjustifiable unless label noise
is a known concern.

\subsection{Runtime Diagnostics}

\begin{table}[h]
  \centering
  \caption{
    Runtime diagnostic reference: healthy operating ranges for EWL monitoring
    quantities.
  }
  \label{tab:diagnostics}
  \begin{tabular}{lll}
    \toprule
    \textbf{Diagnostic} & \textbf{Healthy range} & \textbf{Action if outside} \\
    \midrule
    $\mathcal{H}(w)/\log B$ & $[0.30, 0.85]$ & $<0.30$: increase $\tau$; $>0.85$: decrease $\tau$ \\
    $w_{\max}/w_{\min}$      & $[2\times, 20\times]$ & $>50\times$: possible instability \\
    EMA mean $\bar\mu$       & Decreasing or stable & Rising: check learning rate \\
    Proxy $g^{(t)}$          & $>0$ after warmup     & $\approx 0$ after epoch 5: adapter may be frozen \\
    \bottomrule
  \end{tabular}
\end{table}

\subsection{Limitations}

\paragraph{Marginal clean-data gains.}
On noise-free data, EWL matches SFT within standard deviation across all three
benchmarks ($\Delta \leq +0.16\%$).
The adaptation frontier is not a sufficient proxy for ``difficult but
learnable'' examples when labels are clean, and clean-data performance should
not be cited as a benefit of EWL.

\paragraph{Frequency-blindness.}
EWL's progress signal is purely velocity-based and has no notion of class
frequency.
A rare-class sample that never improves --- because it is genuinely hard, not
noisy --- receives the same near-zero weight as a mislabelled sample.
Combining progress-based weighting with class-aware reweighting is an open
direction.

\paragraph{Single-pass and streaming incompatibility.}
EWL's EMA requires that each training sample is visited multiple times across
epochs.
Single-pass fine-tuning or streaming settings break this assumption --- the
EMA never converges, the signal is unreliable, and the warmup never terminates
meaningfully.
A reservoir or sliding-window approximation is left as future work.

\paragraph{Architecture and task scope.}
All experiments use ViT-Small/16 on fine-grained visual classification.
Whether the adaptation frontier hypothesis generalises to language models,
generation, regression, or other PEFT methods beyond LoRA is not established
here.
